% !PRINTED 136
\weberpage{136}

From (3) and (4) there follows for the cross ratios:
\weq{5}{\frac{(\xi_2\eta_3)(\xi_1\eta_4)}{(\xi_1\eta_3)(\xi_2\eta_4)} = \frac{\vartheta^4_{10}}{\vartheta^4_{00}},}
\weq{6}{\frac{(\xi_1\eta_2)(\xi_3\eta_4)}{(\xi_1\eta_3)(\xi_2\eta_4)} = \frac{\vartheta^4_{01}}{\vartheta^4_{00}}.}
If we now differentiate two of the equations (1), say the first
two, there follows:
\[
\begin{aligned}
2\vartheta_{01}(u)\vartheta'_{01}(u)\,du &= \eta_1\,d\xi - \xi_1\,d\eta = (\eta_1\,d\xi),\\
2\vartheta_{11}(u)\vartheta'_{11}(u)\,du &= \eta_2\,d\xi - \xi_2\,d\eta = (\eta_2\,d\xi),
\end{aligned}
\]
and from this, with the use of (1),
\[
\begin{aligned}
2\vartheta_{01}(u)\vartheta_{11}(u)\bigl[\vartheta'_{01}(u)\vartheta_{11}(u) - \vartheta'_{11}(u)\vartheta_{01}(u)\bigr]\,du \\
= (\eta_1\,d\xi)\vartheta^2_{11}(u) - (\eta_2\,d\xi)\vartheta^2_{01}(u) \\
= (\eta_1\,d\xi)(\xi\eta_2) - (\eta_2\,d\xi)(\xi\eta_1) = (\xi_2\eta_1)(\xi\,d\eta).
\end{aligned}
\]
The last expression can be derived without calculation by regarding the
next to last one as a linear function of \(d\xi\), \(d\eta\) which vanishes for
\(d\xi : d\eta = \xi : \eta\) and is therefore divisible by \((\xi\,d\eta)\). The
quotient \((\xi_2\eta_1)\) is obtained by setting \(d\xi : d\eta = \xi_2 : \eta_2\). With
the use of \S\,23, (6) one then obtains
\weq{7}{2\pi\vartheta^2_{01}\vartheta_{00}(u)\vartheta_{11}(u)\vartheta_{10}(u)\vartheta_{01}(u)\,du = (\xi_1\eta_2)(\xi\,d\eta).}
If one introduces here, according to (3) (first and last formula),
\weq{8}{(\xi_1\eta_2)\vartheta^2_{00} = \sqrt{(\xi_1\eta_3)(\xi_2\eta_4)}\ \vartheta^2_{01}}
and substitutes for the \(\vartheta\)-functions the expressions (1), there
finally follows
\weq{9}{2\pi\vartheta^2_{00}\,du = \frac{\sqrt{(\xi_1\eta_3)(\xi_2\eta_4)}\,(\xi\,d\eta)}{\sqrt{(\xi\eta_1)(\xi\eta_2)(\xi\eta_3)(\xi\eta_4)}},}
by which \(du\) is represented as an \gesperrt{elliptic differential of the first kind}
in homogeneous variables [\S\,1, (5)].

By (1) the ratio \(\xi : \eta\) is determined as a \gesperrt{doubly periodic function}
of \(u\). But likewise the ratios of the square roots
\weq{10}{\sqrt{(\xi\eta_1)},\qquad \sqrt{(\xi\eta_2)},\qquad \sqrt{(\xi\eta_3)},\qquad \sqrt{(\xi\eta_4)}}
are also defined as single-valued doubly periodic functions, and the sign
of the square roots in (9) is thereby, and by (8), likewise uniquely
determined.
